% Thaqafa-RepE: Representation Engineering for Arab Cultural Concepts
%
% Paper scaffold. Every section carries an outline of what belongs in it and,
% where relevant, the exact claim the corresponding code already supports.
% Placeholder numbers are written as \todo{...} so they are impossible to
% mistake for results. Search for "\todo" before submitting anything.
%
% Style: this uses only base LaTeX packages so that it compiles anywhere,
% including CI, without downloading a conference style file. To switch to the
% official ACL template, drop acl.sty / acl_natbib.bst into this directory and
% replace the preamble below with \documentclass[11pt]{article} +
% \usepackage[review]{acl}. Section structure and citations carry over
% unchanged. See README.md in this directory.

\documentclass[11pt,a4paper]{article}

\usepackage[utf8]{inputenc}
\usepackage{amsmath}
\usepackage{amssymb}
\usepackage{booktabs}
\usepackage{graphicx}
\usepackage{natbib}
\usepackage{url}
% Deliberately omitted: fontenc/lmodern/microtype. They improve typography but
% need texlive-fonts-recommended, and this scaffold is meant to compile on a
% bare texlive-latex-base install. The ACL and NeurIPS style files pull in
% their own font setup, so switching to either supersedes this choice.
\usepackage[hidelinks]{hyperref}
\usepackage[margin=1in]{geometry}
\usepackage{xcolor}

% Author notes. Renders in draft mode, silent in final mode:
% \renewcommand{\todo}[1]{}
\newcommand{\todo}[1]{\textcolor{red}{[TODO: #1]}}

% Arabic terms are transliterated in the body text. Rendering Arabic script
% properly needs XeLaTeX or LuaLaTeX with a suitable font (see build.sh --xe);
% under pdflatex the transliteration is what appears.
\newcommand{\arabterm}[2]{\emph{#1} (#2)}

\title{Thaqafa-RepE: Representation Engineering for Arab Cultural Concepts
       in Large Language Models}

\author{
  Thaqafa-RepE Contributors \\
  \texttt{https://github.com/al3obdi/Thaqafa-RepE}
}

\date{\today}

\begin{document}
\maketitle

\begin{abstract}
Alignment research has largely been written in English, by and for a narrow
slice of the world. When a language model reasons about
\arabterm{wasta}{influence through personal connections},
\arabterm{muru'ah}{moral integrity and noble character}, or the obligations of
\arabterm{diyafa}{Arab hospitality}, it tends to fall back on translated
approximations that flatten each concept into its nearest Western analogue and
lose the social weight that makes the original meaningful.

We present Thaqafa-RepE, a reproducible toolkit for studying whether Arab
cultural concepts are linearly represented inside open-weight language models,
and whether those representations can be steered. We contribute (i) a bilingual
dataset of Arab cultural concepts with Arabic and English exemplars and
cultural context annotations; (ii) an implementation of contrastive
mean-difference vector extraction over the residual stream, with a generated
neutral baseline that isolates concept-specific directions from generic
sentence structure; (iii) a hook-based steering mechanism with a scoped context
manager that guarantees the model returns to its unmodified state; and (iv) an
evaluation protocol that reports steering \emph{effect} and fluency \emph{cost}
jointly, and compares representation steering against prompt-engineering
baselines.

\todo{Fill in headline findings once experiments are run: which layers carry
each concept, at what strength steering becomes effective, where fluency breaks
down, and whether steering beats the persona prompting baseline.}
\end{abstract}

%--------------------------------------------------------------------------
\section{Introduction}
\label{sec:introduction}

\todo{Write the full introduction. Structure and the claims to make:}

\begin{itemize}
  \item \textbf{The gap.} Alignment work is overwhelmingly anglophone. Models
        are evaluated on values benchmarks written in English and grounded in
        Western moral intuitions. Concepts central to other cultures are
        handled, if at all, through translation.
  \item \textbf{Why translation is not enough.} \arabterm{wasta}{} is not
        simply "nepotism": it is contested, situationally legitimate, and
        bound up with kinship obligation. A model that maps it onto
        "corruption" has not understood it; it has replaced it. Cite empirical
        evidence of this flattening \citep{naous2024having}.
  \item \textbf{Why representation engineering.} Behavioural probing through
        prompts confounds what a model \emph{represents} with what it is
        willing to \emph{say}. Reading the residual stream directly separates
        the two \citep{zou2023representation}.
  \item \textbf{Research questions.} Legibility: are these concepts linearly
        represented? Locality: at which layers, and is the representation
        shared between Arabic and English descriptions? Controllability: can
        injecting a concept direction make answers more culturally grounded
        without degrading fluency or factuality?
  \item \textbf{Contributions.} Enumerate the four listed in the abstract, with
        a pointer to the released code and dataset.
\end{itemize}

%--------------------------------------------------------------------------
\section{Related Work}
\label{sec:related}

\subsection{Representation Engineering and Activation Steering}
\todo{Position this work against the RepE literature.}
\begin{itemize}
  \item Reading vectors and linear concept directions in the residual stream
        \citep{zou2023representation}.
  \item Activation addition as a steering primitive
        \citep{turner2023activation}; the relationship between steering
        strength and output coherence.
  \item Linear probing as evidence for linearly encoded features
        \citep{alain2017understanding}, and the well-known caveat that a probe
        can succeed without the feature being causally used.
  \item \textbf{Our position:} we treat probe accuracy as evidence of
        \emph{legibility} and steering effect as evidence of \emph{causal use},
        and report them separately rather than conflating them.
\end{itemize}

\subsection{Cultural Bias and Value Alignment in LLMs}
\todo{Survey what is known about cultural coverage in current models.}
\begin{itemize}
  \item Evidence that models default to Western cultural framings even when
        prompted in another language \citep{naous2024having}.
  \item Value and moral benchmarks, and their geographic skew.
  \item Arabic-centric models and evaluation resources, e.g.\ Jais
        \citep{sengupta2023jais}.
\end{itemize}

\subsection{Cross-lingual Representation}
\todo{Connect to the shared-representation literature.}
\begin{itemize}
  \item Evidence for language-agnostic concept representations in multilingual
        models, and the layers at which they emerge.
  \item \textbf{Why it matters here:} if the Arabic and English exemplars of
        one concept yield the same direction, cultural steering transfers
        across the language of the prompt. If they do not, the concept is
        represented twice, and that is a finding in itself
        (\S\ref{sec:results}).
\end{itemize}

%--------------------------------------------------------------------------
\section{Methodology}
\label{sec:methodology}

\subsection{Contrastive Vector Extraction}
\label{sec:extraction}

Let $h^{(\ell)}(x) \in \mathbb{R}^{d}$ denote the residual stream activation
after transformer block $\ell$, averaged over the non-padding tokens of input
$x$. Given a set $P$ of prompts expressing a cultural concept and a set $N$ of
culturally neutral prompts, the concept vector is the normalised difference of
the two means:

\begin{equation}
  \label{eq:extraction}
  v^{(\ell)} =
  \frac{
    \frac{1}{|P|}\sum_{x \in P} h^{(\ell)}(x)
    - \frac{1}{|N|}\sum_{x \in N} h^{(\ell)}(x)
  }{
    \left\lVert
      \frac{1}{|P|}\sum_{x \in P} h^{(\ell)}(x)
      - \frac{1}{|N|}\sum_{x \in N} h^{(\ell)}(x)
    \right\rVert_{2}
  }.
\end{equation}

Subtracting the neutral mean is what makes the direction concept-specific.
Without it, the average is dominated by components shared by every sentence in
the language, and the vectors for distinct concepts point in nearly the same
direction. Three implementation choices matter for reproducibility:

\begin{itemize}
  \item \textbf{Padding is masked.} Batching prompts of unequal length pads the
        short ones; including pad positions in the mean pulls short prompts
        toward the pad embedding. Only leading or trailing pads are masked,
        since GPT-2-style tokenizers reuse a single id for pad, bos and eos.
  \item \textbf{Per-prompt then per-set averaging.} Each prompt is averaged
        over its own real tokens before prompts are averaged together, so a
        long prompt does not outvote a short one.
  \item \textbf{Accumulation in \texttt{float32}} even when the model runs in
        \texttt{bfloat16}.
\end{itemize}

\todo{State how many curated negatives were used per concept, or that the
generated neutral bank was used, and report the sensitivity of the extracted
direction to that choice.}

\subsection{Hook-Based Steering}
\label{sec:steering}

Injection is the mirror image of extraction: a forward hook writes to the same
residual stream point that extraction reads from. For every sequence position
$t$ and every layer $\ell \in L$ in the injection set,

\begin{equation}
  \label{eq:injection}
  h^{(\ell)}_{t} \leftarrow h^{(\ell)}_{t} + \alpha \, v^{(\ell)} .
\end{equation}

Hooks are attached with TransformerLens \citep{nanda2022transformerlens} and
removed through a scoped context manager, so the model provably returns to its
unmodified state after every measurement.

Because $v^{(\ell)}$ has unit norm, the coefficient $\alpha$ is measured in
residual stream norms and is the only magnitude parameter: $\alpha > 0$
amplifies the concept, $\alpha < 0$ suppresses it, and $\alpha = 0$ recovers the
unmodified model exactly. Since a zero offset is a mathematical no-op, the
$\alpha = 0$ condition is measured through the identical code path as the
steered conditions and serves as an exact control.

\subsection{Linear Probes}
\label{sec:probes}

Equation~\ref{eq:extraction} requires choosing $\ell$. Rather than fixing it by
convention, we fit an $L_2$-regularised logistic regression on standardised
activations at each layer, discriminating concept prompts from neutral prompts,
and report $k$-fold cross-validated accuracy against the majority-class floor.
Standardisation matters because residual stream norms grow with depth; an
unstandardised probe would find later layers easier for reasons unrelated to
the concept. Cross-validation matters because with $d \gg n$ a linear
classifier separates almost any labelling of the training set, including a
random one.

\todo{Report $n$ per concept, the number of folds, and the regularisation
strength actually used.}

%--------------------------------------------------------------------------
\section{Experimental Setup}
\label{sec:setup}

\subsection{Models}
\todo{Report the exact checkpoints, revisions, dtype and hardware.}
\begin{itemize}
  \item \texttt{meta-llama/Meta-Llama-3-8B-Instruct} --- the anglophone
        instruction-tuned baseline.
  \item \texttt{inceptionai/jais-13b-chat} --- an Arabic-centric model
        \citep{sengupta2023jais}, to test whether cultural legibility depends
        on Arabic-heavy pretraining.
  \item \todo{Optionally a third model to separate scale from language mix.}
\end{itemize}

\subsection{Dataset}
The concept dataset is distributed with the code as
\texttt{data/datasets/cultural\_concepts.jsonl}, one JSON object per line.
Each entry carries a stable identifier, the concept in Arabic and English, a
category, a one-sentence definition, Arabic and English exemplar sentences,
a cultural-context note, and a sentiment label that is permitted to be
\texttt{mixed} --- contested concepts are recorded as contested rather than
flattered into positives.

\todo{Report the final concept count and per-category breakdown; describe the
native-speaker review protocol and inter-annotator agreement; state dialect
coverage.}

\subsection{Metrics}
\label{sec:metrics}
We report effect and cost separately, because a steering configuration that
leaves fluency untouched because it did nothing must not be mistaken for a
cheap win.

\begin{itemize}
  \item \textbf{Probe accuracy} (\S\ref{sec:probes}): cross-validated, reported
        against the majority-class floor. Evidence of legibility.
  \item \textbf{Steering effect:} the KL divergence between the steered and
        unsteered next-token distributions, averaged over prompts. Evidence
        that the injection did something.
  \item \textbf{Fluency cost:} mean cross-entropy on held-out text the steering
        did not produce. Measured this way, a rise indicates damage to language
        modelling rather than a change of topic.
  \item \textbf{Human evaluation:} blind pairwise ratings of cultural
        grounding, appropriateness and factuality by native speakers.
        \todo{Report rater count, dialect backgrounds, instructions, pay, and
        agreement statistics. The automatic metrics detect damage; only this
        can establish success.}
\end{itemize}

\subsection{Baselines}
\label{sec:baselines}
Representation steering must be compared against the intervention any
practitioner would try first: asking the model in plain language. We evaluate a
no-instruction control, a direct instruction, and a persona framing
(\S\ref{sec:results}). The two interventions are not symmetric --- steering
perturbs the model, prompting only changes the input --- so cross-entropy on the
\emph{prompt} is not comparable across them. We instead score the fluency of
each condition's \emph{output} under the same unmodified model.

%--------------------------------------------------------------------------
\section{Results}
\label{sec:results}

\subsection{Where are cultural concepts linearly readable?}
\todo{Insert the layer sweep figure and describe the shape of the curve.}

\begin{figure}[t]
  \centering
  % \includegraphics[width=\columnwidth]{figures/layer_sweep.pdf}
  \fbox{\parbox[c][3.5cm][c]{0.9\linewidth}{\centering
    \todo{Figure: probe accuracy against layer index, one line per concept,
    with the majority-class floor marked.}}}
  \caption{Cross-validated probe accuracy by layer. Accuracy above the dashed
  chance line indicates the concept is linearly readable at that depth.}
  \label{fig:layer-sweep}
\end{figure}

\subsection{Steering effect against fluency cost}
\todo{Insert the strength sweep and the layer-set grid.}

\begin{table}[t]
  \centering
  \begin{tabular}{lrrr}
    \toprule
    Layers & $\alpha$ & Effect (KL) & Cost ($\Delta$ CE) \\
    \midrule
    $\{\ell^{*}\}$            & 1.0 & \todo{} & \todo{} \\
    $\{\ell^{*}\}$            & 2.0 & \todo{} & \todo{} \\
    $\{\ell^{*}-1,\ell^{*}\}$ & 1.0 & \todo{} & \todo{} \\
    \bottomrule
  \end{tabular}
  \caption{Effect and cost for each injection configuration. Spreading the same
  direction across more layers is expected to raise both.}
  \label{tab:layer-sets}
\end{table}

\subsection{Steering versus prompting}
\todo{Report the head-to-head comparison from \S\ref{sec:baselines}.}

\begin{table}[t]
  \centering
  \begin{tabular}{lrrr}
    \toprule
    Condition & Grounding (human) & Fluency & Extra tokens \\
    \midrule
    No instruction     & \todo{} & \todo{} & 0 \\
    Direct instruction & \todo{} & \todo{} & \todo{} \\
    Persona framing    & \todo{} & \todo{} & \todo{} \\
    Steering ($\alpha=\todo{}$) & \todo{} & \todo{} & 0 \\
    \bottomrule
  \end{tabular}
  \caption{Representation steering against prompt-engineering baselines. The
  persona framing is the hardest baseline and the one the claim rests on.}
  \label{tab:baselines}
\end{table}

\subsection{Do Arabic and English exemplars share a direction?}
\todo{Report the cosine similarity between vectors extracted from the Arabic-only
and English-only exemplars of each concept, and whether steering with one
transfers to prompts in the other.}

\subsection{Does suppression remove the concept?}
\todo{Negative $\alpha$ should remove the concept, not merely produce generic
text. Report whether suppressed outputs remain on-topic.}

%--------------------------------------------------------------------------
\section{Limitations}
\label{sec:limitations}

\todo{State these plainly; reviewers will find them regardless.}
\begin{itemize}
  \item \textbf{Dataset scale and provenance.} A small concept inventory,
        authored by a limited set of annotators, cannot represent "Arab
        culture" as a whole; dialect and regional coverage is uneven.
  \item \textbf{Probes are correlational.} High probe accuracy shows a
        direction is readable, not that the model uses it. The steering
        experiments are what speak to causal use.
  \item \textbf{Automatic metrics detect damage, not success.} Cultural
        grounding is established only by the human evaluation.
  \item \textbf{Essentialism risk.} Encoding a culture as a set of vectors
        risks reifying contested practices into fixed model behaviour.
        \arabterm{wasta}{} in particular is normatively disputed within the
        societies where it operates, which is why it carries a \texttt{mixed}
        sentiment label.
  \item \textbf{Dual use.} The same mechanism that amplifies a cultural concept
        can suppress it.
\end{itemize}

%--------------------------------------------------------------------------
\section{Conclusion}
\label{sec:conclusion}

\todo{Write once the results are in. The shape of the argument:}
\begin{itemize}
  \item Restate what was measured and what was found about legibility,
        locality and controllability.
  \item Say plainly whether steering beat the prompting baseline, and at what
        cost in fluency.
  \item \textbf{Future work:} curated minimal-pair negatives; more concepts and
        wider dialect coverage; extending the method to other
        under-represented cultures using the released toolkit; and combining
        steering with prompting rather than treating them as alternatives.
\end{itemize}

%--------------------------------------------------------------------------
\section*{Ethics Statement}
\todo{Required by ACL. Cover: consent and compensation for native-speaker
annotators; the essentialism risk noted in \S\ref{sec:limitations}; the dual-use
concern that suppression is as easy as amplification; and the decision to label
contested concepts as contested.}

\section*{Reproducibility Statement}
All code, the concept dataset, and the exact commands used to produce every
number in this paper are released at
\url{https://github.com/al3obdi/Thaqafa-RepE}.
\todo{Add the release commit hash, the resolved dependency lockfile, random
seeds, and total compute.}

%--------------------------------------------------------------------------
\bibliographystyle{plainnat}
\bibliography{references}

\end{document}
